\documentclass{article}

\usepackage{amsmath,amssymb,amsthm}

\newtheorem{theorem}{Theorem}

\begin{document}

\begin{theorem}
For $0<\alpha<1$, let $p_\alpha(j)$ the pmf of sibuya distribution:
$$
p_\alpha(j) :=\frac{\alpha\prod_{i=1}^{j-1}(i-\alpha)}{j!} = \alpha \frac{\Gamma(j-\alpha)}{\Gamma(j+1)\Gamma(1-\alpha)},
\qquad j\ge 1.
$$
Then it holds that 
$$
\mathrm{fisher}(\alpha)
:=
\alpha^{-2}
+
\sum_{j=1}^\infty p_\alpha(j)\frac{1}{\alpha(j-\alpha)}
=
\frac{\pi}{\alpha\sin(\pi\alpha)}.
$$
Using $\sin(\pi\alpha) < \pi(1-\alpha)$, we know that the fisher information is larger than $\frac{1}{\alpha(1-\alpha)}$. 
\end{theorem}


\begin{proof}
Using
$$
\frac{1}{j-\alpha}=\int_0^1 x^{j-\alpha-1}\,dx,
\qquad j\ge 1,
$$
we obtain
$$
\sum_{j=1}^\infty \frac{p_\alpha(j)}{j-\alpha}
=
\int_0^1 x^{-\alpha-1}\sum_{j=1}^\infty p_\alpha(j)x^j\,dx.
$$
Now the probability generating function of the Sibuya distribution is
$$
\sum_{j=1}^\infty p_\alpha(j)x^j
=
1-(1-x)^\alpha,
\qquad x\in[0,1].
$$
Therefore,
$$
\sum_{j=1}^\infty \frac{p_\alpha(j)}{j-\alpha}
=
\int_0^1 x^{-\alpha-1}\bigl(1-(1-x)^\alpha\bigr)\,dx.
$$
Hence
$$
\mathrm{fisher}(\alpha)
=
\alpha^{-2}
+
\alpha^{-1}\int_0^1 x^{-\alpha-1}\bigl(1-(1-x)^\alpha\bigr)\,dx.
$$
Using integration by parts, 
\begin{align*}
\int_0^1 x^{-\alpha-1}\bigl(1-(1-x)^\alpha\bigr)\,dx
&=
\left[
-\frac{1}{\alpha}x^{-\alpha}\bigl(1-(1-x)^\alpha\bigr)
\right]_0^1
+
\int_0^1 x^{-\alpha}(1-x)^{\alpha-1}\,dx \\
&=
-\frac{1}{\alpha}
+
B(1-\alpha,\alpha).
\end{align*}
where $B$ is the Beta function. 
Since
$$
B(1-\alpha,\alpha)=\Gamma(1-\alpha)\Gamma(\alpha),
$$
it follows that
$$
\mathrm{fisher}(\alpha)
=
\alpha^{-2}
+
\alpha^{-1}\left(
-\alpha^{-1}+\Gamma(1-\alpha)\Gamma(\alpha)
\right)
=
\frac{\Gamma(1-\alpha)\Gamma(\alpha)}{\alpha}.
$$
Finally, Euler's reflection formula gives
$$
\Gamma(\alpha)\Gamma(1-\alpha)=\frac{\pi}{\sin(\pi\alpha)},
$$
and thus
$$
\mathrm{fisher}(\alpha)
=
\frac{\pi}{\alpha\sin(\pi\alpha)}.
$$
\end{proof}

\newpage

Let $k_{n,j}$ be the emperical frequency generated by the Gibbs partition and let $P_{n,j}=\frac{k_{n,j}}{k_n}$. We know that the pointwise convergence hold
$$
P_{n,j}\to^p p_\alpha(j).
$$
As a corollary, one can show
\begin{align}\label{eq:bounded}
    \sum_{j} |P_{n,j} - p_{\alpha}(j)| \to 0. 
\end{align}

We aim to show the asymptotic normality of the QMLE under a certain assumption on the distributional convergence of the random measure $\sqrt{K_n} (P_{n,\cdot}-p_{\alpha})$ on $\mathbb{N}$. 
Recall that the QMLE is defined as the root of the system 
$$
\Psi_n(x) = 0
$$
where $\Psi_n$ is the random function given by 
\begin{align*}
    \Psi_n(x) &= \frac{1}{k_n} \Bigl(\frac{k_n - 1}{x} - \sum_{j=2}^{n} k_{n,j} \sum_{i=1}^{j-1} \frac{1}{i - x}\Bigr) = -\frac{1}{k_n} + \frac{1}{x} - \sum_{j=2}^n P_{n,j} \sum_{i=1}^{j-1} \frac{1}{i - x}
\end{align*}
($\Psi_n$ is the derivative of the log likelihood of the Ewens--Pitman with $\theta=0$ normalized by $k_n$). We also define the deterministic function $\Psi$ as 
\begin{align*}
\Psi(x) &= \frac{1}{\alpha} - \sum_{j=2}^\infty p_{\alpha}(j) \frac{1}{i-x}.
\end{align*}
This is obtained from $\Psi_n$ with $P_{n,j}$ replaced by $p_\alpha(j)$ and ignoring the diminishing term $1/k_n$. We showed in previous paper that the map 
$\Psi:(0,1)\to\mathbb{R}$ satisfies 
\begin{itemize}
    \item $\Psi:(0,1)\to\mathbb{R}$ is $C^1$. 
    \item $\Psi$ is strictly decreasing and 
    $$   
    \Psi'(x)=-\frac{1}{x^2} - \sum_{j=2}^\infty p_\alpha(j) \sum_{i=1}^{j-1}\frac{1}{(i-x)^2}<0
    $$
    \item $\Psi'(\alpha)=-{i}(\alpha) < 0$, where ${i}(\alpha)$ is the Fisher Information of the discrete distribution $p_\alpha(j)$. 
    \item $\Psi$ has its unique root at $x=\alpha$, i.e., $\Psi(\alpha)=0$.
    \item $\Psi_n'(x)$
\end{itemize}
Furthermore, using \eqref{eq:bounded}, one can show that for any closed interval $I=[s,t]$ with $0<s<t<1$, 
\begin{align*}
 \sup_{x\in I}|\Psi_n(x)'- \Psi(x)'|\to^p 0. 
\end{align*}

Now, using $\hat{\Psi}_n(\hat{\alpha}_n)=\Psi(\alpha)=0$, by Taylor's expansion, 
\begin{align*}
    \Psi_n(\alpha) - \Psi(\alpha) = \Psi_n(\alpha)-\Psi_n(\hat{\alpha}_n) = \Psi_n'(\tilde{\alpha}_n) (\alpha-\hat\alpha_n).
\end{align*}


Now we assume that for certain test function $f:\mathbb{N}\to\mathbb{R}$, 
$$
\sqrt{K_n} \Bigl(\sum_j f(j) (P_{n,j} - p_\alpha(j))\Bigr) \to N(0, f^\top \Gamma_\alpha f) 
$$
where $\Gamma_\alpha = D(p_{\alpha}) - p_{\alpha} p_{\alpha}^\top$


Now, using your $L_2$ distributional convergence, taking test function $f_x:\mathbb{N}\to \mathbb{R}$ as $f_x(j) = \sum_{i=1}^{j-1}\frac{1}{i-x}$ for each $x$, we would get 
\begin{align*}
   \forall x\in (0,1), \quad  \sqrt{K_n} (\Psi_n(x) - \Psi(x)) \to N(0, h_x^\top \Gamma_\alpha h_x)
\end{align*}
This in particular $\Psi_n(x)\to^p \Psi(x)$ pointwisely and hence we have the consistency of the QMLE. 

Inportantly, taking $x=\alpha$.  
By the same argument, now taking the test function $h_x:\mathbb{N}\to \mathbb{R}$ as $h_x(j) = \sum_{i=1}^{j-1}\frac{1}{(i-x)^2}$ for each $x$, we also get
\begin{align*}
   \forall x\in (0,1), \quad  \Psi_n'(x) \to^p \Psi'(x). 
\end{align*}


\end{document}